\documentclass[aspectratio=169,11pt]{beamer}

% -----------------------------------------------------------------------------
% Procesamiento Masivo de Datos
% Linux + shell + pipes + regex + sed + awk
% Reorganized for a progressive, hands-on introduction.
% -----------------------------------------------------------------------------
\usepackage[T1]{fontenc}
\usepackage[utf8]{inputenc}
\usepackage[spanish,es-nodecimaldot]{babel}
\usepackage{microtype}
\usepackage{booktabs}
\usepackage{array}
\usepackage{listings}
\usepackage{tikz}
\usepackage{xurl}
\usetikzlibrary{arrows.meta,positioning,shapes.geometric,calc,overlay-beamer-styles}

\definecolor{UOHBlue}{RGB}{0,117,201}
\usetheme[accent=UOHBlue]{Jacquenetta}
\usepackage{uoh-jacquenetta}

% --- Course metadata ---------------------------------------------------------
\title[Procesamiento Masivo de Datos]{Procesamiento Masivo de Datos}
\subtitle{Linux, tuberías, expresiones regulares, \texttt{sed} y \texttt{awk}}
\author{Alex Di Genova}
\institute{Universidad de O'Higgins}
\date{2026}

\begin{document}

\titleframe

% =============================================================================
\section{Motivación}
% =============================================================================

\begin{frame}[fragile]{La idea central del curso}
  \centering
  \vspace{3mm}
  \metric{Blue}{Datos}{GB \textrightarrow{} TB \textrightarrow{} PB}\hfill
  \metric{Mint}{Herramientas}{pequeñas y componibles}\hfill
  \metric{Coral}{Objetivo}{flujos reproducibles}
  \vspace{8mm}
  \begin{wideblock}{Pregunta guía}
    ¿Cómo procesamos archivos grandes \alert{sin abrirlos completos en memoria}, evitando pasos manuales y dejando un análisis que pueda repetirse?
  \end{wideblock}
  \vfill
  \small\color{Slate}En esta clase, la terminal no es el objetivo: es el lenguaje común para mover, inspeccionar, transformar y resumir datos.
\end{frame}

\begin{frame}[fragile]{Objetivos de aprendizaje}
  \begin{columns}[T,onlytextwidth]
    \column{.48\textwidth}
      \begin{wideblock}{Linux y shell}
        \begin{itemize}
          \item navegar y manipular archivos
          \item comprender rutas, permisos y procesos
          \item usar redirecciones y tuberías
        \end{itemize}
      \end{wideblock}
    \column{.48\textwidth}
      \begin{wideblock}{Procesamiento de texto}
        \begin{itemize}
          \item buscar patrones con regex y \term{grep}
          \item transformar flujos con \term{sed}
          \item filtrar, calcular y resumir con \term{awk}
        \end{itemize}
      \end{wideblock}
  \end{columns}
  \vspace{5mm}
  \centering\pill{Meta de la clase}\quad Convertir una pregunta sobre datos en una tubería verificable y reutilizable.
\end{frame}

\begin{frame}[fragile]{Del archivo al resultado: un modelo mental}
  \centering
  \begin{tikzpicture}[>=Latex,node distance=3mm]
    \tikzset{stage/.style={rounded corners=4pt,minimum width=23mm,minimum height=12mm,align=center,font=\small\bfseries}}
    \node[stage,fill=Blue!16,visible on=<1->] (f) {archivo\\\scriptsize entrada};
    \node[stage,fill=Cyan!18,right=of f,visible on=<2->] (i) {inspeccionar\\\scriptsize head · wc · file};
    \node[stage,fill=Amber!20,right=of i,visible on=<3->] (s) {seleccionar\\\scriptsize grep · cut};
    \node[stage,fill=Mint!18,right=of s,visible on=<4->] (t) {transformar\\\scriptsize sed · awk};
    \node[stage,fill=Coral!18,right=of t,visible on=<5->] (o) {resultado\\\scriptsize tabla · resumen};
    \draw[->,thick,Blue,visible on=<2->] (f)--(i);
    \draw[->,thick,Blue,visible on=<3->] (i)--(s);
    \draw[->,thick,Blue,visible on=<4->] (s)--(t);
    \draw[->,thick,Blue,visible on=<5->] (t)--(o);
  \end{tikzpicture}
  \vspace{8mm}
  \begin{wideblock}{Principio Unix}
    Resolver problemas complejos combinando herramientas simples, cada una especializada en una tarea pequeña.
  \end{wideblock}
\end{frame}

% =============================================================================
\section{Linux y shell: lo esencial}
% =============================================================================

\begin{frame}[fragile]{¿Qué ocurre cuando escribimos un comando?}
  \centering
  \begin{tikzpicture}[node distance=.8mm]
    \tikzset{layer/.style={minimum width=118mm,minimum height=8mm,rounded corners=2pt,font=\small\bfseries}}
    \node[layer,fill=Blue!16,visible on=<1->] (u) {usuario: terminal · shell · programas};
    \node[layer,fill=Cyan!18,below=of u,visible on=<2->] (s) {llamadas al sistema};
    \node[layer,fill=Mint!18,below=of s,visible on=<3->] (k) {kernel: procesos · memoria · red · sistema de archivos};
    \node[layer,fill=Ink,text=Paper,below=of k,visible on=<4->] (h) {CPU · RAM · discos · interfaces de red};
  \end{tikzpicture}
  \vspace{6mm}
  \begin{wideblock}{Idea clave}
    La \alert{shell} interpreta el comando y crea procesos; el \alert{kernel} administra los recursos físicos.
  \end{wideblock}
\end{frame}

\begin{frame}[fragile]{La shell es interfaz y también lenguaje}
  \begin{columns}[T,onlytextwidth]
    \column{.43\textwidth}
      \textbf{Como interfaz}
      \begin{itemize}
        \item ejecuta programas
        \item conecta entrada y salida
        \item controla procesos
      \end{itemize}
      \textbf{Como lenguaje}
      \begin{itemize}
        \item variables y condiciones
        \item ciclos y funciones
        \item scripts versionables
      \end{itemize}
    \column{.54\textwidth}
\begin{lstlisting}
$ pwd
/home/alex/project
$ command --option input.tsv
\end{lstlisting}
      \small\pill{programa} \term{command}\quad
      \pill[Cyan]{opción} \term{--option}\quad
      \pill[Mint]{argumento} \term{input.tsv}
  \end{columns}
\end{frame}

\begin{frame}[fragile]{Navegar: saber dónde estoy y qué hay alrededor}
\codebar{shell}
\begin{lstlisting}
pwd                    # directorio actual
ls                     # listar
ls -lh                  # tamanos legibles
ls -la                  # incluye archivos ocultos
cd data                 # entrar a data/
cd ..                   # subir un nivel
cd ~                    # volver al HOME
\end{lstlisting}
  \vspace{3mm}
  \centering
  \term{.} = aquí \qquad \term{..} = directorio padre \qquad \term{\~} = directorio personal
\end{frame}

\begin{frame}[fragile]{Archivos y directorios: crear, copiar, mover, borrar}
  \begin{columns}[T,onlytextwidth]
    \column{.48\textwidth}
\codebar{crear y copiar}
\begin{lstlisting}
mkdir results
mkdir -p results/qc
cp input.tsv backup.tsv
cp -r data/ data_backup/
\end{lstlisting}
    \column{.48\textwidth}
\codebar{mover y borrar}
\begin{lstlisting}
mv old.tsv new.tsv
mv new.tsv results/
rm file.tmp
rm -r directory/
\end{lstlisting}
  \end{columns}
  \vspace{5mm}
  \pill[Coral]{Atención}\quad \term{rm} no tiene papelera. Revise rutas y comodines antes de ejecutar.
\end{frame}

\begin{frame}[fragile]{Encontrar y medir datos}
\begin{lstlisting}
find data/ -type f -name '*.tsv'
find . -type f -size +1G

du -sh data/
du -h --max-depth=1 data/ | sort -h

df -h
\end{lstlisting}
  \vspace{4mm}
  \begin{wideblock}{Dos preguntas distintas}
    \term{du} pregunta: ``¿cuánto ocupa este directorio?'' \qquad
    \term{df} pregunta: ``¿cuánto espacio queda en el sistema de archivos?''
  \end{wideblock}
\end{frame}

\begin{frame}[fragile]{Comodines y comillas: la shell expande antes de ejecutar}
  \begin{columns}[T,onlytextwidth]
    \column{.47\textwidth}
      \textbf{Globbing}
\begin{lstlisting}
ls *.tsv
ls sample?.txt
ls run_[123].log
\end{lstlisting}
    \column{.49\textwidth}
      \textbf{Comillas}
\begin{lstlisting}
echo "$HOME"
echo '$HOME'
rm -- "file with spaces.txt"
\end{lstlisting}
  \end{columns}
  \vspace{6mm}
  \begin{wideblock}{Regla práctica}
    Use comillas dobles alrededor de variables y rutas cuando no quiera que espacios o comodines cambien el significado del comando.
  \end{wideblock}
\end{frame}

\begin{frame}[fragile]{Pedir ayuda es parte del trabajo}
\begin{lstlisting}
man sort
sort --help
command -v awk
type cd
\end{lstlisting}
  \vspace{5mm}
  \begin{columns}[T,onlytextwidth]
    \column{.31\textwidth}\pill{man}\quad manual completo
    \column{.31\textwidth}\pill[Cyan]{--help}\quad recordatorio rápido
    \column{.31\textwidth}\pill[Mint]{command -v}\quad ¿qué ejecutable usaré?
  \end{columns}
\end{frame}

\begin{frame}[fragile]{Permisos: quién puede hacer qué}
\begin{lstlisting}
$ ls -l analysis.sh
-rwxr-x--- 1 alex lab 2048 Aug 31 analysis.sh
\end{lstlisting}
  \begin{center}
  \begin{tabular}{>{\bfseries}lccc}
    \toprule
    & usuario & grupo & otros \\
    \midrule
    representación & rwx & r-x & --- \\
    significado & leer/escribir/ejecutar & leer/ejecutar & ninguno \\
    \bottomrule
  \end{tabular}
  \end{center}
  \vspace{3mm}
  \pill{mínimo privilegio}\quad Conceda sólo el acceso necesario; evite \term{chmod 777} como solución rápida.
\end{frame}

% =============================================================================
\section{Flujos, redirecciones y tuberías}
% =============================================================================

\begin{frame}[fragile]{Tres flujos; infinitas combinaciones}
  \centering
  \begin{tikzpicture}[>=Latex,node distance=12mm]
    \node[draw=Blue,very thick,rounded corners=4pt,minimum width=35mm,minimum height=18mm,font=\bfseries] (cmd) {comando};
    \node[left=of cmd,align=center] (in) {\pill[Cyan]{stdin · 0}\\datos de entrada};
    \node[right=of cmd,yshift=7mm,align=center] (out) {\pill[Mint]{stdout · 1}\\resultado};
    \node[right=of cmd,yshift=-10mm,align=center] (err) {\pill[Coral]{stderr · 2}\\diagnóstico};
    \draw[->,line width=1.2pt,Cyan] (in)--(cmd);
    \draw[->,line width=1.2pt,Mint] (cmd)--(out);
    \draw[->,line width=1.2pt,Coral] (cmd)--(err);
  \end{tikzpicture}
  \vspace{5mm}
\begin{lstlisting}
command < input.txt > result.txt 2> errors.log
\end{lstlisting}
\end{frame}

\begin{frame}[fragile]{Redirecciones: decidir de dónde vienen y adónde van los datos}
\begin{lstlisting}
command > result.txt        # sobrescribe stdout
command >> result.txt       # agrega stdout al final
command 2> errors.log       # guarda stderr
command > out.txt 2>&1      # stdout + stderr
command < input.txt         # toma stdin desde archivo
\end{lstlisting}
  \vspace{4mm}
  \pill[Amber]{Consejo}\quad Antes de usar \term{>}, confirme que no está sobrescribiendo un resultado importante.
\end{frame}

\begin{frame}[fragile]{La tubería: stdout de un proceso \textrightarrow{} stdin del siguiente}
\codebar{pipeline.sh}
\begin{lstlisting}
cat file.txt \
  | grep 'ERROR' \
  | sort \
  | uniq -c \
  | sort -nr
\end{lstlisting}
  \vspace{4mm}
  \begin{wideblock}{Regla de depuración}
    Construya la tubería de izquierda a derecha. Ejecute cada etapa sobre una muestra pequeña antes de agregar la siguiente.
  \end{wideblock}
\end{frame}

\begin{frame}[fragile]{No siempre necesitamos \texttt{cat}}
  \begin{columns}[T,onlytextwidth]
    \column{.48\textwidth}
      \textbf{Funciona}
\begin{lstlisting}
cat file.tsv | grep chile
\end{lstlisting}
    \column{.48\textwidth}
      \textbf{Más directo}
\begin{lstlisting}
grep chile file.tsv
\end{lstlisting}
  \end{columns}
  \vspace{6mm}
  \small
  \textbf{Pero} \term{cat} es útil para concatenar archivos, alimentar comandos que sólo leen stdin o hacer una tubería más explícita durante la enseñanza.
\end{frame}

% =============================================================================
\section{Inspeccionar, seleccionar y ordenar}
% =============================================================================

\begin{frame}[fragile]{Inspeccionar antes de transformar}
\begin{lstlisting}
head -n 5 worldcitiespop.csv
tail -n 5 worldcitiespop.csv
wc -l worldcitiespop.csv
file worldcitiespop.csv
less worldcitiespop.csv
\end{lstlisting}
  \vspace{4mm}
  \begin{columns}[T,onlytextwidth]
    \column{.31\textwidth}\pill{1}\quad ¿hay cabecera?
    \column{.31\textwidth}\pill{2}\quad ¿cuál es el delimitador?
    \column{.31\textwidth}\pill{3}\quad ¿qué valores faltan?
  \end{columns}
  \vspace{5mm}
  \alert{Nunca empiece transformando un archivo que todavía no ha inspeccionado.}
\end{frame}

\begin{frame}[fragile]{Una caja de herramientas para texto}
  \begin{center}
  \renewcommand{\arraystretch}{1.25}
  \begin{tabular}{>{\bfseries\color{Blue}}l p{.69\textwidth}}
    \toprule
    inspeccionar & \term{head}, \term{tail}, \term{less}, \term{wc}, \term{file} \\
    buscar & \term{grep} \\
    columnas & \term{cut}, \term{paste} \\
    ordenar/agrupar & \term{sort}, \term{uniq} \\
    transformar & \term{tr}, \term{sed}, \term{awk} \\
    combinar/comparar & \term{join}, \term{comm}, \term{diff} \\
    encontrar archivos & \term{find} \\
    \bottomrule
  \end{tabular}
  \end{center}
  \vspace{3mm}
  \centering\small Elija la herramienta más simple que exprese correctamente la operación.
\end{frame}

\begin{frame}[fragile]{Seleccionar columnas con \texttt{cut}}
\codebar{worldcitiespop.csv}
\begin{lstlisting}
head -n 3 worldcitiespop.csv

cut -d, -f1,2,5 worldcitiespop.csv | head
\end{lstlisting}
  \vspace{4mm}
  \begin{wideblock}{Importante}
    \term{cut -d,} sirve aquí porque el archivo es simple. Un CSV general puede contener comas entre comillas y necesita un parser real.
  \end{wideblock}
\end{frame}

\begin{frame}[fragile]{Ordenar y contar categorías}
\begin{lstlisting}
cut -d, -f1 worldcitiespop.csv \
  | tail -n +2 \
  | sort \
  | uniq -c \
  | sort -nr \
  | head
\end{lstlisting}
  \vspace{4mm}
  \centering
  \pill{sort}\; agrupa iguales \quad + \quad
  \pill[Cyan]{uniq -c}\; cuenta adyacentes \quad + \quad
  \pill[Mint]{sort -nr}\; ordena por frecuencia
\end{frame}

\begin{frame}[fragile]{Caso guiado: las 10 ciudades más pobladas}
  \textbf{Pregunta:} ¿cuáles son las 10 ciudades con mayor población?
\codebar{top-cities.sh}
\begin{lstlisting}
tail -n +2 worldcitiespop.csv \
  | sort -t, -k5,5nr \
  | head -n 10
\end{lstlisting}
  \begin{center}
  \begin{tikzpicture}[>=Latex,node distance=4mm]
    \node[rounded corners=3pt,fill=Blue!16,minimum width=34mm,minimum height=9mm,font=\scriptsize\bfseries] (n0) {quitar cabecera};
    \node[rounded corners=3pt,fill=Amber!18,minimum width=38mm,minimum height=9mm,font=\scriptsize\bfseries,right=of n0] (n1) {ordenar campo 5 numérico};
    \node[rounded corners=3pt,fill=Mint!18,minimum width=34mm,minimum height=9mm,font=\scriptsize\bfseries,right=of n1] (n2) {tomar diez};
    \draw[->,thick] (n0)--(n1); \draw[->,thick] (n1)--(n2);
  \end{tikzpicture}
  \end{center}
\end{frame}

\begin{frame}[fragile]{Procesar datos comprimidos sin descomprimirlos a disco}
\begin{lstlisting}
gzip -dc reads.txt.gz | head
zcat table.tsv.gz | wc -l

zcat table.tsv.gz \
  | grep -v '^#' \
  | cut -f1,2 \
  | head
\end{lstlisting}
  \vspace{4mm}
  \begin{wideblock}{Por qué importa}
    En datos masivos, evitar copias intermedias ahorra espacio, I/O y tiempo. Muchas herramientas pueden leer directamente desde \term{stdin}.
  \end{wideblock}
\end{frame}

\begin{frame}[fragile]{¿Qué operaciones escalan bien?}
  \begin{columns}[T,onlytextwidth]
    \column{.47\textwidth}
      \begin{wideblock}{Streaming}
        \term{grep}, \term{cut}, \term{sed} y muchos usos de \term{awk} procesan una línea a la vez.
      \end{wideblock}
    \column{.47\textwidth}
      \begin{wideblock}{Requieren más recursos}
        \term{sort} necesita ordenar globalmente; agregaciones en \term{awk} pueden almacenar muchas claves en memoria.
      \end{wideblock}
  \end{columns}
  \vspace{6mm}
  \centering\pill[Amber]{Pregunta de escalabilidad}\quad ¿cuánto crece la memoria y el disco temporal cuando crece el archivo?
\end{frame}

% =============================================================================
\section{Expresiones regulares y grep}
% =============================================================================

\begin{frame}[fragile]{Regex: describir patrones, no posiciones fijas}
  \begin{center}
  \begin{tabular}{l p{.67\textwidth}}
    \toprule
    \textbf{Patrón} & \textbf{Significado} \\
    \midrule
    \term{\char94} \quad \term{\$} & inicio y final de línea \\
    \term{.} & cualquier carácter \\
    \term{[abc]} & uno de los caracteres del conjunto \\
    \term{[\char94 abc]} & cualquier carácter excepto los del conjunto \\
    \term{*} \quad \term{+} \quad \term{?} & cero o más; uno o más; opcional \\
    \term{\{m,n\}} & entre \emph{m} y \emph{n} repeticiones \\
    \term{(a|b)} & agrupación y alternativas en ERE \\
    \term{[[:digit:]]} & clase portable de dígitos \\
    \bottomrule
  \end{tabular}
  \end{center}
  \vspace{4mm}
  \pill[Amber]{Consejo}\quad Use \term{grep -E} y \term{sed -E} para expresiones regulares extendidas.
\end{frame}

\begin{frame}[fragile]{Buscar con \texttt{grep}}
\begin{lstlisting}
grep 'Chile' countries.txt
grep -i 'chile' countries.txt
grep -n 'ERROR' pipeline.log
grep -v '^#' variants.tsv
grep -E '^(chr1|chr2)[[:space:]]' variants.tsv
\end{lstlisting}
  \vspace{4mm}
  \begin{columns}[T,onlytextwidth]
    \column{.23\textwidth}\pill{i}\; ignora caso
    \column{.23\textwidth}\pill[Cyan]{n}\; muestra línea
    \column{.23\textwidth}\pill[Mint]{v}\; invierte
    \column{.23\textwidth}\pill[Coral]{E}\; regex extendida
  \end{columns}
\end{frame}

\begin{frame}[fragile]{Regex en acción}
\codebar{regex-demo.sh}
\begin{lstlisting}
$ printf '%s\n' ac abc abbc abbbc abbbbbc \
    | grep -E '^ab{1,3}c$'
abc
abbc
abbbc
\end{lstlisting}
  \vspace{3mm}
  \begin{wideblock}{Leer el patrón de izquierda a derecha}
    inicio · \term{a} · entre 1 y 3 letras \term{b} · \term{c} · final de línea
  \end{wideblock}
\end{frame}

% =============================================================================
\section{sed: transformar un flujo}
% =============================================================================

\begin{frame}[fragile]{\texttt{sed}: pensar en ``buscar y transformar''}
  \begin{columns}[T,onlytextwidth]
    \column{.48\textwidth}
      \textbf{Sustituir}
\codebar{sed · sustitución}
\begin{lstlisting}
sed 's/hello/bye/' file.txt
sed 's/hello/bye/g' file.txt
\end{lstlisting}
    \column{.48\textwidth}
      \textbf{Normalizar espacios}
\codebar{sed · regex extendida}
\begin{lstlisting}
sed -E 's/[[:space:]]+/ /g' file.txt
\end{lstlisting}
  \end{columns}
  \vspace{5mm}
  \centering \term{s/patrón/reemplazo/banderas}
\end{frame}

\begin{frame}[fragile]{\texttt{sed}: seleccionar líneas y rangos}
\begin{lstlisting}
sed -n '10,20p' file.txt
sed -n '/ERROR/p' pipeline.log
sed '10,20s/hello/bye/g' file.txt
sed '/^#/d' variants.tsv
\end{lstlisting}
  \vspace{4mm}
  \begin{columns}[T,onlytextwidth]
    \column{.31\textwidth}\pill{p}\quad imprimir
    \column{.31\textwidth}\pill[Cyan]{d}\quad eliminar
    \column{.31\textwidth}\pill[Mint]{10,20}\quad rango de líneas
  \end{columns}
\end{frame}

\begin{frame}[fragile]{Capturas: reorganizar texto con grupos}
\codebar{sed · grupos de captura}
\begin{lstlisting}
$ echo 'chr1:100-250' \
    | sed -E 's/([^:]+):([0-9]+)-([0-9]+)/\1\t\2\t\3/'
chr1    100    250
\end{lstlisting}
  \vspace{5mm}
  \begin{wideblock}{Idea clave}
    Los paréntesis capturan partes del patrón; \term{\textbackslash1}, \term{\textbackslash2}, \term{\textbackslash3} las reutilizan en el reemplazo.
  \end{wideblock}
\end{frame}

\begin{frame}[fragile]{Ejemplos de transformación con \texttt{sed}}
\begin{lstlisting}
sed 's/2022/2024/g' file.txt
sed 's/2022/2024/g; s/casa/home/g' file.txt

echo 'cats dog bee parrot foxed' \
  | sed -E 's/cat|dog|fox/--/g'
# --s -- bee parrot --ed

echo 'car bat cod map scat dot abacus' \
  | sed 's/b.*at/-/'
# car - dot abacus
\end{lstlisting}
  \vspace{2mm}
  \small\color{Slate}Discusión: \term{.*} es codicioso; consume todo lo posible antes de permitir que el resto del patrón coincida.
\end{frame}

% =============================================================================
\section{awk: filtrar, calcular y resumir}
% =============================================================================

\begin{frame}[fragile]{\texttt{awk}: datos como registros y campos}
  \begin{columns}[T,onlytextwidth]
    \column{.46\textwidth}
      \begin{itemize}
        \item un registro por línea
        \item campos \term{\$1 ... \$NF}
        \item \term{NR}: número de registro
        \item regla = patrón + acción
      \end{itemize}
    \column{.51\textwidth}
\codebar{awk · filtro}
\begin{lstlisting}
awk -F, 'NR > 1 && $5 > 1e6 {
  print $2, $5
}' worldcitiespop.csv
\end{lstlisting}
  \end{columns}
  \vspace{5mm}
  \centering\term{awk 'condición \{ acción \}' archivo}
\end{frame}

\begin{frame}[fragile]{Separadores: \texttt{FS} para leer, \texttt{OFS} para escribir}
\begin{lstlisting}
awk -F, 'BEGIN {OFS="\t"}
  NR > 1 {print $1, $2, $5}
' worldcitiespop.csv | head
\end{lstlisting}
  \vspace{5mm}
  \begin{wideblock}{Modelo}
    \term{FS} define cómo se separan los campos de entrada; \term{OFS} define cómo se separan los campos impresos.
  \end{wideblock}
\end{frame}

\begin{frame}[fragile]{Patrones y condiciones con \texttt{awk}}
\begin{lstlisting}
awk '$NF < 0' table.txt
awk 'NR == 1 || $5 > 1000000' worldcitiespop.csv
awk '$2 ~ /error/' log.txt
awk 'tolower($2) ~ /error/' log.txt
awk 'NF >= 5 {print $1, $NF}' table.txt
\end{lstlisting}
  \vspace{4mm}
  \pill{\textasciitilde}\; coincide con regex \quad
  \pill[Cyan]{!\textasciitilde}\; no coincide \quad
  \pill[Mint]{NF}\; número de campos
\end{frame}

\begin{frame}[fragile]{El ciclo de vida de \texttt{awk}}
  \centering
  \begin{tikzpicture}[>=Latex,node distance=9mm]
    \tikzset{phase/.style={rounded corners=4pt,minimum width=31mm,minimum height=14mm,align=center,font=\bfseries}}
    \node[phase,fill=Blue!16] (b) {BEGIN\\\scriptsize inicializar};
    \node[phase,fill=Cyan!18,right=of b] (m) {cada registro\\\scriptsize filtrar · calcular};
    \node[phase,fill=Mint!18,right=of m] (e) {END\\\scriptsize resumir};
    \draw[->,thick,Blue] (b)--(m); \draw[->,thick,Blue] (m)--(e);
    \draw[->,thick,Slate,loop above] (m) to node[above,font=\scriptsize]{siguiente línea} (m);
  \end{tikzpicture}
\begin{lstlisting}
awk 'BEGIN {sum=0}
     {sum += $1}
     END {print sum}' values.txt
\end{lstlisting}
\end{frame}

\begin{frame}[fragile]{Agregación con arreglos asociativos}
  \textbf{Pregunta:} ¿cuál es la población total registrada por país?
\codebar{population-by-country.sh}
\begin{lstlisting}
awk -F, '
  NR > 1 && $5 != "" { pop[$1] += $5 }
  END { for (c in pop) print c, pop[c] }
' worldcitiespop.csv \
  | sort -k2,2nr \
  | head
\end{lstlisting}
  \vspace{3mm}
  \pill{clave}\; país \quad + \quad \pill[Cyan]{valor}\; suma acumulada \quad = \quad \pill[Mint]{resumen}\; por grupo
\end{frame}

\begin{frame}[fragile]{Contar categorías con \texttt{awk}}
\begin{lstlisting}
awk -F, '
  NR > 1 { count[$1]++ }
  END {
    for (c in count)
      print c, count[c]
  }
' worldcitiespop.csv \
  | sort -k2,2nr \
  | head
\end{lstlisting}
  \vspace{4mm}
  \small\color{Slate}El patrón \term{count[key]++} aparece constantemente en logs, tablas, variantes, muestras y resultados de pipelines.
\end{frame}

\begin{frame}[fragile]{Funciones útiles de \texttt{awk}}
  \begin{columns}[T,onlytextwidth]
    \column{.47\textwidth}
      \textbf{Cadenas}
      \begin{itemize}
        \item \term{length(s)}, \term{substr(s,p,n)}
        \item \term{split(s,a,sep)}
        \item \term{sub}, \term{gsub}, \term{match}
        \item \term{tolower}, \term{toupper}
      \end{itemize}
    \column{.47\textwidth}
      \textbf{Cálculo}
      \begin{itemize}
        \item \term{+ - * / \% \^}
        \item \term{sqrt}, \term{log}, \term{int}
        \item \term{== != < <= > >=}
        \item \term{\&\& || !}
      \end{itemize}
  \end{columns}
\end{frame}

\begin{frame}[fragile]{Laboratorio corto: selección con \texttt{awk}}
\begin{lstlisting}[basicstyle=\ttfamily\fontsize{7}{8}\selectfont\color{Paper}]
$ cat table.txt
brown bread mat hair 42
blue cake mug shirt -7
yellow banana window shoes 3.14

$ awk '{print $2}' table.txt
bread
cake
banana

$ awk '$NF < 0' table.txt
blue cake mug shirt -7

$ awk '{gsub(/b/, "B", $1)} 1' table.txt
Brown bread mat hair 42
Blue cake mug shirt -7
yellow banana window shoes 3.14
\end{lstlisting}
\end{frame}

% =============================================================================
\section{Combinar herramientas}
% =============================================================================

\begin{frame}[fragile]{Ejemplo integrado: frecuencia de palabras}
\codebar{word-frequency.sh}
\begin{lstlisting}
tr '[:upper:]' '[:lower:]' < text.txt \
  | tr -cs '[:alnum:]' '\n' \
  | sed '/^$/d' \
  | sort \
  | uniq -c \
  | sort -nr \
  | head -n 20
\end{lstlisting}
  \vspace{3mm}
  \textbf{Discuta:} ¿qué ocurre con tildes, guiones, números y distintas configuraciones regionales?
\end{frame}

\begin{frame}[fragile]{Ejemplo integrado: resumir un archivo de variantes}
  \small Supongamos un TSV simplificado: \term{CHROM POS REF ALT SAMPLE}
\codebar{variants.tsv}
\begin{lstlisting}
grep -v '^#' variants.tsv \
  | awk 'BEGIN {OFS="\t"}
         {count[$1]++}
         END {for (chr in count) print chr, count[chr]}' \
  | sort -k2,2nr
\end{lstlisting}
  \vspace{4mm}
  \begin{wideblock}{Transferencia de aprendizaje}
    El mismo patrón sirve para contar eventos por cromosoma, errores por nodo, muestras por grupo o registros por categoría.
  \end{wideblock}
\end{frame}

\begin{frame}[fragile]{Combinar dos tablas con \texttt{join}}
  \small \term{join} requiere que ambas tablas estén ordenadas por la clave.
\begin{lstlisting}
sort -k1,1 countries.tsv > countries.sorted.tsv
sort -k1,1 population.tsv > population.sorted.tsv

join -t $'\t' -1 1 -2 1 \
  countries.sorted.tsv population.sorted.tsv
\end{lstlisting}
  \vspace{4mm}
  \pill[Amber]{Concepto}\quad Es el equivalente en shell de un \emph{inner join} simple.
\end{frame}

\begin{frame}[fragile]{Una tubería reproducible deja huellas}
  \begin{columns}[T,onlytextwidth]
    \column{.48\textwidth}
      \begin{itemize}[<+->]
        \item script versionado
        \item entradas identificadas
        \item parámetros explícitos
        \item versiones de herramientas
      \end{itemize}
    \column{.48\textwidth}
      \begin{itemize}[<+->]
        \item errores registrados
        \item conteos de control
        \item salidas deterministas
        \item comando de ejecución documentado
      \end{itemize}
  \end{columns}
  \vspace{5mm}
\begin{lstlisting}
#!/usr/bin/env bash
set -euo pipefail

pipeline input.tsv > result.tsv 2> pipeline.log
\end{lstlisting}
\end{frame}

\begin{frame}[fragile]{Ver el flujo mientras se ejecuta: \texttt{tee}}
\begin{lstlisting}
grep -v '^#' variants.tsv \
  | tee filtered.tsv \
  | awk '{count[$1]++}
         END {for (c in count) print c, count[c]}'
\end{lstlisting}
  \vspace{5mm}
  \begin{wideblock}{Útil para depurar}
    \term{tee} copia el flujo a un archivo y, al mismo tiempo, lo deja continuar por la tubería.
  \end{wideblock}
\end{frame}

% =============================================================================
\section{Trabajo remoto y clúster}
% =============================================================================

\begin{frame}[fragile]{El clúster: dónde ejecutaremos los análisis grandes}
  \begin{columns}[c,onlytextwidth]
    \column{.63\textwidth}
    \centering
    \begin{tikzpicture}[>=Latex,node distance=7mm]
      \tikzset{box/.style={rounded corners=3pt,minimum width=27mm,minimum height=11mm,align=center,font=\small\bfseries}}
      \node[box,fill=Blue!18] (login) {nodo de acceso};
      \node[box,fill=Amber!22,below=of login] (sched) {planificador\\SLURM};
      \node[box,fill=Mint!20,right=16mm of sched] (fs) {almacenamiento\\compartido};
      \node[box,fill=Cyan!18,below left=8mm and -5mm of sched] (n1) {cómputo 01};
      \node[box,fill=Cyan!18,below=8mm of sched] (n2) {cómputo 02};
      \node[box,fill=Cyan!18,below right=8mm and -5mm of sched] (n3) {cómputo n};
      \draw[->,thick,Blue] (login)--node[right,font=\scriptsize]{envía jobs}(sched);
      \foreach \n in {n1,n2,n3}{\draw[->,thick,Blue] (sched)--(\n);\draw[-,thick,Mint] (fs)--(\n);}
    \end{tikzpicture}
    \column{.34\textwidth}
      \begin{itemize}[<+->]
        \item SSH para entrar
        \item SLURM para computar
        \item filesystem compartido
        \item nodos para trabajo pesado
      \end{itemize}
      \vspace{2mm}\alert{No ejecute análisis pesado en el nodo de acceso.}
  \end{columns}
\end{frame}

\begin{frame}[fragile]{Acceso remoto y transferencia}
  \begin{columns}[T,onlytextwidth]
    \column{.48\textwidth}
      \textbf{Conectar}
\begin{lstlisting}
ssh user@cluster
ssh -p 8765 user@cluster
\end{lstlisting}
    \column{.48\textwidth}
      \textbf{Transferir}
\begin{lstlisting}
scp file.tsv user@cluster:/data/
rsync -avP results/ \
  user@cluster:/data/results/
\end{lstlisting}
  \end{columns}
  \vspace{4mm}
  \pill[Mint]{rsync}\quad especialmente útil para directorios grandes y transferencias que pueden reanudarse.
\end{frame}

\begin{frame}[fragile]{Sesiones persistentes para trabajo interactivo}
\begin{lstlisting}
tmux new -s analysis

# Ctrl-b d   -> separar la sesion

tmux ls
tmux attach -t analysis
\end{lstlisting}
  \vspace{4mm}
  \begin{wideblock}{En un clúster}
    Use \term{tmux} para sesiones interactivas livianas; use el planificador para análisis que consumen CPU o memoria de manera significativa.
  \end{wideblock}
\end{frame}

% =============================================================================
\section{Práctica}
% =============================================================================

\begin{frame}[fragile]{Laboratorio: ciudades del mundo}
\codebar{worldcitiespop.csv}
\begin{lstlisting}[basicstyle=\ttfamily\scriptsize\color{Paper}]
head worldcitiespop.csv
wc -l worldcitiespop.csv
head -100 worldcitiespop.csv > sample_worldcitiespop.csv

cut -d, -f1,5 worldcitiespop.csv > country_population.csv

tail -n +2 worldcitiespop.csv \
  | sort -t, -k5,5nr \
  | head -10

awk -F, 'NR > 1 && $5 > 1000000' worldcitiespop.csv
\end{lstlisting}
\end{frame}

\begin{frame}[fragile]{Laboratorio: de pregunta a solución}
  \begin{enumerate}
    \item ¿Cuántas filas tiene el archivo y cuántas poseen población informada?
    \item ¿Cuántos países distintos aparecen?
    \item ¿Cuáles son las 10 ciudades más pobladas?
    \item ¿Cuánta población total registra cada país?
    \item ¿Qué países tienen más ciudades con población $>1$ millón?
  \end{enumerate}
  \vspace{4mm}
  \begin{wideblock}{Criterio de evaluación}
    Cada respuesta debe incluir: comando, resultado, una comprobación y una explicación de una línea.
  \end{wideblock}
\end{frame}

\begin{frame}[fragile]{Desafíos: de simple a compuesto}
  \begin{enumerate}
    \item Contar la frecuencia de palabras en un texto.
    \item Reemplazar \term{hello} por \term{bye} sólo entre las líneas 10 y 20.
    \item Extraer un rango \term{chr:start-end} y convertirlo a tres columnas.
    \item Combinar dos tablas por un campo común.
    \item Contar registros por categoría con un arreglo asociativo de \term{awk}.
    \item Procesar un archivo \term{.gz} sin crear una copia descomprimida.
  \end{enumerate}
  \vspace{3mm}
  \pill[Amber]{Criterio}\quad Primero una prueba pequeña; después el archivo completo.
\end{frame}

% =============================================================================
\section{Herramientas y frontera práctica}
% =============================================================================

\begin{frame}[fragile]{¿Cuándo dejar shell, sed o awk?}
  \begin{columns}[T,onlytextwidth]
    \column{.31\textwidth}
      \pill{Shell}\par\smallskip
      \small orquestar programas, archivos y procesos.
    \column{.31\textwidth}
      \pill[Cyan]{sed / awk}\par\smallskip
      \small transformar y resumir texto estructurado de forma compacta.
    \column{.31\textwidth}
      \pill[Mint]{Python / R}\par\smallskip
      \small lógica extensa, parsers complejos, modelos y visualización.
  \end{columns}
  \vspace{8mm}
  \begin{wideblock}{Frontera práctica}
    Una tubería debe ser fácil de leer. Si la lógica deja de ser obvia, muévala a un script más expresivo.
  \end{wideblock}
\end{frame}

\begin{frame}[fragile]{CSV y tablas modernas}
  \begin{columns}[T,onlytextwidth]
    \column{.31\textwidth}
      \pill{csvtk}\par\smallskip
      \small filtrar, mutar y combinar CSV/TSV dentro de tuberías.
    \column{.31\textwidth}
      \pill[Cyan]{xsv}\par\smallskip
      \small seleccionar, indexar, unir y analizar CSV rápidamente.
    \column{.31\textwidth}
      \pill[Mint]{YouPlot}\par\smallskip
      \small inspeccionar distribuciones directamente en terminal.
  \end{columns}
  \vspace{8mm}
  \centering\small Shell sigue siendo la capa de composición; las herramientas especializadas resuelven formatos que requieren un parser real.
\end{frame}

\begin{frame}[fragile]{Recursos del curso}
  \small
  \begin{itemize}
    \item Linux Upskill Challenge: \url{https://linuxupskillchallenge.org/}
    \item GNU awk: \url{https://www.gnu.org/software/gawk/manual/}
    \item GNU sed: \url{https://www.gnu.org/software/sed/manual/}
    \item Learn GNU awk: \url{https://learnbyexample.github.io/learn_gnuawk/}
    \item Learn GNU sed: \url{https://learnbyexample.github.io/learn_gnused/}
    \item Regex101: \url{https://regex101.com/}
    \item worldcitiespop.csv: \url{https://burntsushi.net/stuff/worldcitiespop.csv}
    \item csvtk: \url{https://github.com/shenwei356/csvtk}
  \end{itemize}
\end{frame}

\begin{frame}[fragile]{Ideas para llevar}
  \begin{enumerate}
    \item Inspeccione antes de transformar.
    \item Piense en datos como flujos: \term{stdin}, \term{stdout}, \term{stderr}.
    \item Combine herramientas pequeñas con tuberías.
    \item \term{grep} busca; \term{sed} transforma; \term{awk} filtra, calcula y resume.
    \item Pregunte siempre cómo escalan memoria, disco e I/O.
    \item Convierta el análisis final en un script reproducible.
  \end{enumerate}
  \vspace{6mm}
  \begin{wideblock}{Siguiente paso}
    Tome una tarea manual, exprésela como tubería, pruébela con una muestra pequeña y luego ejecútela sobre el conjunto completo.
  \end{wideblock}
\end{frame}

\thanksframe[Preguntas]

\end{document}
